\chapter{Product Planning}

\section{Two-Session Overview}

Product planning converts the earlier opportunity work into a clear development direction. In these two sessions, your group will study the market, identify relevant technologies and constraints, and then convert that analysis into a mission statement that can guide later work on customer needs and product specifications.

The exercises for Product Planning is divided into two parts: In the first session, each member will conduct an individual market analysis and draft a mission statement. In the second session, the group will consolidate those analyses into one shared plan and one final mission statement.

By the end of Sessions 5 and 6, your group should be able to:
\begin{enumerate}
    \item identify relevant market segments, stakeholders, technologies, and constraints for the selected hydroponic product;
    \item explain how technical and physics-related considerations influence the planning decisions;
    \item draft and refine a mission statement for the project; and
    \item produce a shared planning document that will support Sessions 7 and 8.
\end{enumerate}

\section{Key Terms}

The following terms are used throughout this chapter:
\begin{itemize}
    \item \textbf{Market segment}: a group of customers with similar needs, behaviours, or use contexts.
    \item \textbf{Technology integration}: the choice of technical features or subsystems that could strengthen the product.
    \item \textbf{Mission statement}: a concise planning statement describing the product, its benefits, target market, and key goals.
    \item \textbf{Stakeholders}: people or groups affected by the product, including users, manufacturers, maintainers, sellers, and investors.
    \item \textbf{Assumptions and constraints}: the working conditions, limits, and uncertainties that shape the project.
\end{itemize}

\section{Why Product Planning Matters}

A good plan prevents later sessions from becoming disconnected. For example, if your team expects the product to be affordable for small urban growers, that assumption will affect customer needs, material choices, manufacturing approach, and target price. Similarly, if the concept depends on sensors or automated flow control, the team must already recognise the implications for calibration, power use, durability, and maintenance.

\section{Session 5: Individual Market Analysis and Draft Mission Statement}

Session 5 focuses on preparation. Each member should review the opportunity from earlier sessions and gather evidence that will help the team plan realistically.

\subsection{Activity 1: Market Analysis}

Before class, analyse the planning questions below for your selected product.



To prepare for product planning, address the following questions:

\begin{enumerate}

\item \textbf{Market Segments}
Identify and list potential market segments that should be considered when designing the product and its features. For example, are you designing for household hobby growers, educational users, or commercial operators?

\item \textbf{Technology Integration}
Identify any new or emerging technologies that could be integrated into the product to enhance functionality, efficiency, or user experience. Examples include sensors, modular structures, dosing systems, or low-power control units.

\item \textbf{Manufacturing and Service}
Define the key manufacturing and service objectives, including any technical, logistical, or operational constraints that may influence product development.

\item \textbf{Target Market and Price}
\begin{enumerate}
    \item Clearly define the primary and secondary target market(s) for the product.
    \item Establish a realistic target price range in Malaysian Ringgit (RM),for the product within these markets.
\end{enumerate}

\item \textbf{Assumptions, Constraints, and Stakeholders}
\begin{enumerate}
    \item State the key assumptions underlying your product plan.
    \item Identify constraints that may limit product development, production, or market entry.
    \item List the main stakeholders involved in the product planning process.
\end{enumerate}

\end{enumerate}


For Industrial Physics students, this analysis should also consider physical behaviour where relevant. For example, flow stability, corrosion resistance, thermal exposure, or sensor accuracy may affect both technology selection and product feasibility.

\subsection{Market Analysis Template}

Use Table~\ref{tab:session56-market-analysis} to organise your findings.

\begin{table}[H]
\centering
\caption{Example market analysis template for hydroponic product planning}
\label{tab:session56-market-analysis}
\renewcommand{\arraystretch}{1.2}
\begin{tabularx}{\textwidth}{|>{\raggedright\arraybackslash}p{0.17\textwidth}|>{\raggedright\arraybackslash}p{0.18\textwidth}|>{\raggedright\arraybackslash}X|>{\raggedright\arraybackslash}p{0.16\textwidth}|}
\hline
\textbf{Planning element} & \textbf{Example entry} & \textbf{Questions to answer} & \textbf{Sources or notes} \\
\hline
Market segment & Urban home grower & Who uses the product and what problem do they face? &  \\
\hline
Technology opportunity & Low-cost water-level sensor & What technology could improve function, safety, or usability? &  \\
\hline
Manufacturing or service issue & Easy cleaning of wetted parts & What production or maintenance constraint matters most? &  \\
\hline
Target market and price & Beginner hobby growers, RM150--RM250 & What price range is realistic and why? &  \\
\hline
Assumptions and constraints & Indoor use, limited power, small footprint & What limits must the concept respect? &  \\
\hline
Stakeholders & User, retailer, fabricator, maintainer & Who affects or is affected by the product? &  \\
\hline
\end{tabularx}
\end{table}

\subsection{Activity 2: Draft the Mission Statement}

After completing the analysis, write a short draft mission statement. A useful mission statement should include:


\begin{enumerate}

\item \textbf{Product Description}  
Provide a concise description of the product.

\item \textbf{Benefit Proposition}  
Clearly describe the key benefit(s) the product offers to its target market(s).

\item \textbf{Key Business Goals}  
State the primary business goals the product is intended to achieve (e.g., market share, revenue growth, profitability).

\item \textbf{Primary Market}  
Identify the main target market segment for the product.

\item \textbf{Secondary Markets}  
Identify any additional market segments that the product may also serve.

\item \textbf{Assumptions and Constraints}  
Briefly list the key assumptions and constraints identified in Activity~1 that are relevant to the mission statement.

\item \textbf{Stakeholders}  
List the primary stakeholders who are relevant to the product’s success.

\end{enumerate}




\subsection{Mission Statement Template}

You may draft the mission statement using the structure in Table~\ref{tab:session56-mission-template} before converting it into a short paragraph.

\begin{table}[htbp]
\centering
\caption{Mission statement drafting template}
\label{tab:session56-mission-template}
\renewcommand{\arraystretch}{1.2}
\begin{tabularx}{\textwidth}{|>{\raggedright\arraybackslash}p{0.22\textwidth}|X|}
\hline
\textbf{Field} & \textbf{Draft content} \\
\hline
Product description &  \\
\hline
Benefit proposition &  \\
\hline
Primary market &  \\
\hline
Secondary market &  \\
\hline
Key goal &  \\
\hline
Assumptions and constraints &  \\
\hline
Stakeholders &  \\
\hline
\end{tabularx}
\end{table}

\subsection{Expected Session 5 Output}

Before the next class, each member should bring:
\begin{itemize}
    \item a completed market analysis table;
    \item one draft mission statement; and
    \item brief notes on supporting sources or evidence.
\end{itemize}

\section{Consolidating the Group Plan}

Session 6 focuses on synthesis. Your group will compare the individual analyses, resolve differences, and produce one shared planning direction.

\subsection{Activity 1: Merge the Market Analyses}

Work through the following steps as a team:
\begin{enumerate}
    \item compare the market segments identified by each member;
    \item agree on the most relevant customer group or groups;
    \item discuss which technologies are attractive and technically realistic;
    \item review manufacturing, service, and maintenance issues that the product must respect;
    \item confirm a realistic target market and price range; and
    \item finalise the most important assumptions, constraints, and stakeholders.
\end{enumerate}

If opinions differ, use the innovation charter, user value, and project feasibility as the main criteria for deciding.

\subsection{Group Consolidation Template}

Use Table~\ref{tab:session56-group-consolidation} to capture the final consensus.

\begin{table}[htbp]
\centering
\caption{Group market analysis consolidation template}
\label{tab:session56-group-consolidation}
\renewcommand{\arraystretch}{1.2}
\begin{tabularx}{\textwidth}{|>{\raggedright\arraybackslash}p{0.17\textwidth}|>{\raggedright\arraybackslash}X|>{\raggedright\arraybackslash}p{0.18\textwidth}|}
\hline
\textbf{Planning element} & \textbf{Consensus decision} & \textbf{Rationale or evidence} \\
\hline
Market segment &  &  \\
\hline
Technology opportunity &  &  \\
\hline
Manufacturing and service &  &  \\
\hline
Target market and price &  &  \\
\hline
Assumptions and constraints &  &  \\
\hline
Stakeholders &  &  \\
\hline
\end{tabularx}
\end{table}

\subsection{Activity 2: Finalise the Mission Statement}

Once the group analysis is consolidated, convert it into one final mission statement. Check the statement against the following questions:
\begin{enumerate}
    \item Is the product described clearly?
    \item Is the value proposition meaningful for the selected market?
    \item Are the main assumptions and constraints realistic?
    \item Does the statement align with the innovation charter and shortlisted opportunity?
    \item Does the wording support later work on customer needs and specifications?
\end{enumerate}

\subsection{Collaboration Guidance}

To keep the discussion fair and efficient:
\begin{itemize}
    \item let each member present their strongest evidence briefly;
    \item record disagreements before voting or choosing a final direction;
    \item verify claims with credible sources where possible; and
    \item make sure the final wording is understood by the full team.
\end{itemize}

\subsection*{Optional Remarks (May or Not yield Extra Credit)}

To demonstrate a deeper understanding and potentially earn extra mark be it for individual or group submission, you can add remarks that provide further justification, context, or insights to your analysis in either the pre-class or in-class activities.

\textbf{Guidelines for Remarks:}

\begin{enumerate}
\item \textbf{Detailed Justification:} If you make a specific choice or assumption in your analysis (e.g., selecting a particular target market, choosing a specific technology), provide a detailed justification for this decision. Explain why this is a suitable choice and support your reasoning.
\item \textbf{Reference/Citation:} Back up your justifications or insights with credible sources. This could include:
\begin{itemize}
\item \textbf{Internet Sources}: Cite reputable websites, industry reports, market research data, or news articles relevant to your product or market. Ensure the sources are credible (e.g., Statista, Gartner, industry-specific publications, academic institutions).
\item \textbf{Journal Articles/Academic Research:} If applicable, cite relevant academic articles or research papers that support your analysis or provide theoretical frameworks. You can use databases like Google Scholar, JSTOR, or IEEE Xplore to find relevant publications.
\item \textbf{Textbook or Course Materials:} Refer to concepts, frameworks, or examples discussed in your textbook or other course materials to strengthen your arguments.
\end{itemize}
\item \textbf{Beyond the Basics:} Remarks should go beyond simply answering the basic questions. They should demonstrate critical thinking, industry awareness, and a deeper engagement with the product planning process. Examples of valuable remarks could include:
\begin{itemize}
\item Discussing potential risks or challenges associated with your chosen market segment or technology.
\item Analyzing the competitive landscape and how it influences your product planning decisions.
\item Exploring relevant industry trends and how they might impact your product's success.
\item Justifying your chosen target price point based on market research or competitor analysis.
\item Explaining the rationale behind prioritizing certain stakeholders over others.
\end{itemize}
\item \textbf{Proper Citation Format:} Use a consistent citation format (i.e, APA) when referencing external sources. Even informal internet sources should be clearly linked or described.
\end{enumerate}

\section{Summary and Next Steps}

At the end of Session 6, your group should have one refined market analysis and one final mission statement. These outputs become the planning foundation for Session 7, where you will identify customer needs, and Session 8, where those needs will be translated into measurable specifications.
